\subsection*{\theappendixitem.\ 局部环的膨胀}
\phantomsection
\addcontentsline{toc}{subsection}{\theappendixitem.\ 局部环的膨胀}

设 A 是一个局部环。

我们用 A{]}X{[} 表示多项式环 A{[}X{]} 在素理想 \(m_{A}A[X]\) 处的局部环。这是一个以 \(m_{A}A{]}X{[}\) 为极大理想的局部环；典范同态 A → A{]}X{[} 是局部且平坦的，且 A{]}X{[} 的剩余域是由 X 的类生成的 \(\kappa_{A}\) 的纯扩张。

引理 1. — 设 \(\mathbf{P} \in \mathbf{A}[\mathbf{X}]\) 是一个首一多项式，其像 \(\overline{\mathbf{P}}\) 在 \(\kappa_{\mathbf{A}}[\mathbf{X}]\) 中不可约。则 A-代数 \(\mathbf{B} = \mathbf{A}[\mathbf{X}] / (\mathbf{P})\) 是局部的且在 \(\mathbf{A}\) 上有限，其极大理想为 \(\mathfrak{m}_{\mathbf{A}}\mathbf{B}\)；典范同态 \(\rho: \mathbf{A} \to \mathbf{B}\) 是局部且平坦的，剩余域扩张 \(\kappa_{\mathbf{A}} \to \kappa_{\mathbf{B}}\) 是代数的且由 \(x\) （\(\mathbf{X}\) 的类）所生成，而 \(x\) 在 \(\kappa_{\mathbf{A}}\) 上的最小多项式是 \(\overline{\mathbf{P}}\) 。

由于多项式 P 是首一的，A-模 B 是有限型的自由模（A，IV，第 10 页）。环 B/m \(_{A}\) B 等同于 \(\kappa_{A}[X]/(\overline{P})\) ，因而是一个域；于是理想 m \(_{A}\) B 是极大的。设 q 是 B 的一个极大理想；则理想 \(\rho^{-1}(q)\) 是极大的（V，§ 2，第 1 目，命题 1）；因此我们有 \(\rho^{-1}(q)=m_{A}\) ，从而 q ⊃ m \(_{A}\) ，最后 q = m \(_{A}\) B。于是环 B 是局部的。由此立即得到引理 1。

定义 1. --- 设 A 是一个局部环。若 A-代数 B 同构于 A-代数 A{]}X{[}，或者存在 A{[}X{]} 中的一个首一多项式 P，其在 \(\kappa_{\mathrm{A}}[\mathrm{X}]\) 中的像不可约，使得 B 同构于 A-代数 A{[}X{]}/(P)，则称 A-代数 B 是 A 的一个初等膨胀。

设 B 是 A 的一个初等膨胀。由前面的讨论可得下列性质：

\begin{enumerate}
\def\labelenumi{\alph{enumi})}
\item
  环 B 是局部的，且从 A 到 B 的典范同态是局部且平坦的，特别还是单射的（I，§ 3，第 5 目，命题 8）。
\item
  B 的剩余域 \(\kappa_{B}\) 是 A 的剩余域 \(\kappa_{A}\) 的一个单生成扩张。若 \(\kappa_{B}\) 在 \(\kappa_{A}\) 上的次数 d 有限，则 B 是秩为 d 的自由 A-模。
\item
  我们有 \(m_{B} = m_{A}B\) 。特别地，若 A 是域，则 B 也是域。一个域扩张是初等膨胀当且仅当它是单生成的。
\item
  若 A 是 Noether 的，则 B 也是 Noether 的。
\end{enumerate}

定义 2. --- 设 A 是一个局部环。若存在一个具有最大元 ω 的良序集 Λ，以及 B 的子代数的一个递增族 (B \(_{λ}\) ) \(_{λ∈Λ}\) ，满足下列条件，则称 A-代数 B 是 A 的一个膨胀：

\begin{enumerate}
\def\labelenumi{\alph{enumi})}
\item
  我们有 \(\mathbf{B}_{\omega} = \mathbf{B}\) ，且环 \(\mathbf{B}_{\lambda}\) 对所有 \(\lambda \in \Lambda\) 都是局部的。
\item
  若 \(\alpha\) 是 \(\Lambda\) 的最小元，则 A-代数 \(\mathbf{B}_{\alpha}\) 同构于 \(\mathbf{A}\) 。
\item
  设 \(v \neq \alpha\) 在 \(\Lambda\) 中，并设 \(S_{v}\) 是由 \(\lambda \in \Lambda\) 中满足 \(\lambda < v\) 者组成的集合。若 \(S_{v}\) 没有最大元，我们有 \(B_{v} = \bigcup_{\lambda \in S_{v}} B_{\lambda}\) ；若 \(S_{v}\) 有最大元 \(\mu\) ，则 \(B_{v}\) 是 \(B_{\mu}\) 的一个初等膨胀。
\end{enumerate}

设 B 是一个环，\(\rho: A \to B\) 是一个环同态。若由 \(\rho\) 定义的 A-代数具有该性质，则称 \(\rho\) 是一个膨胀（相应地，初等膨胀）。若如此，则 \(\rho\) 是单射。

例. --- 1) 每个域扩张都是一个膨胀。事实上，设 K 是域 k 的一个扩张。给 K 赋予一个以 0 为最大元的良序，且对 \(\lambda \in \mathbf{K}\) ，设 \(\mathbf{K}_{\lambda}\) 为 K 的一个子 \(k\) -扩张，它由 K 的元素 \(\beta\) （满足 \(\beta < \lambda\) ）生成。条件 \(a)\) ，\(b)\) ，\(c)\) 对于 \(k\) ，K 以及族 \((\mathbf{K}_{\lambda})_{\lambda \in \mathbf{K}}\) 的验证是直接的。

\begin{enumerate}
\def\labelenumi{\arabic{enumi})}
\setcounter{enumi}{1}
\tightlist
\item
  设 A 是一个局部环，I 是一个指标集。用 A{]} \((X_{i})_{i\in I}\) {[} 表示多项式环 A \([(X_{i})_{i\in I}]\) 在素理想 \(m_{A}A[(X_{i})_{i\in I}]\) 处的局部环。A-代数 A{]} \((X_{i})_{i\in I}\) {[} 是 A 的一个膨胀。事实上，给集合 I 赋予良序；设 Λ 是由 I 添加一个最大元 ω 而得到的良序集。对 i ∈ I，把 A{]} \((X_{j})_{j<i}\) {[} 等同于 B \(_{i}\) ，后者是 B = A{]} \((X_{i})_{i\in I}\) 的一个子代数，并令 B \(_{\omega}\) = B。族 (B \(_{\lambda}\) ) \(_{\lambda\in\Lambda}\) 满足条件 a)，b)，c)。
\end{enumerate}

注记. --- 采用定义 2 的记号，环 \(\mathbf{B}_{\mu}\) 是 \(\mathbf{B}_{\lambda}\) 的一个膨胀，当 \(\lambda \leqslant \mu\) 时。

命题 2. --- 设 A 是一个局部环，B 是 A 的一个膨胀。

\begin{enumerate}
\def\labelenumi{\alph{enumi})}
\item
  环 B 是局部的，且我们有 \(\mathfrak{m}_{\mathrm{A}}\mathbf{B} = \mathfrak{m}_{\mathrm{B}}\) 。
\item
  A-代数 B 是忠实平坦的。
\item
  典范同态
\end{enumerate}

\[
\gamma_ {\mathrm{B}}: \operatorname{gr} (\mathrm{A}) \otimes_ {\kappa_ {\mathrm{A}}} \kappa_ {\mathrm{B}} \rightarrow \operatorname{gr} (\mathrm{B})
\]

是双射。

\begin{enumerate}
\def\labelenumi{\alph{enumi})}
\setcounter{enumi}{3}
\tightlist
\item
  若 A 是 Noether 的，则 B 也是 Noether 的，并且 A 与 B 的 Hilbert-Samuel 级数（VIII，§ 4，第 3 目）相等。
\end{enumerate}

设 \((\mathbf{B}_{\lambda})_{\lambda\in\Lambda}\) 是 B 的子代数的一个族，满足定义 2 的条件 a)，b)，c)。

设 \(\Lambda'\) 是由满足下列性质的指标 \(\lambda \in \Lambda\) 组成的集合：对满足 \(\mu \leqslant \lambda\) 且属于 \(\Lambda\) 的每个 μ，A-代数 \(B_{\mu}\) 是局部且忠实平坦的，且 \(m_{B_{\mu}} = m_A B_{\mu}\) 。假设 \(\Lambda' \neq \Lambda\) ，并设 v 是 \(\Lambda - \Lambda'\) 的最小元。我们有 \(\alpha \in \Lambda'\) ，从而 v \(\neq \alpha\) 。\(S_{v}\) 包含于 \(\Lambda'\) 。若 \(S_{v}\) 没有最大元，我们有 \(B_{v} = \bigcup_{\lambda \in S_{v}} B_{\lambda}\) ，且 v 属于 \(\Lambda'\)

  这由第 1 目的命题 1 得出。若 \(S_v\) 有最大元 \(\mu\) ，则有 \(\mu \in \Lambda'\) 且 \(B_v\) 是 \(B_\mu\) 的一个初等膨胀：由定义 1 后面的注记，我们仍有 \(v \in \Lambda'\) ，由此得到矛盾。

当 A 是 Noether 环时，可以用类似的方式证明，集合 \(\Lambda''\) （由指标 \(\lambda \in \Lambda\) 中使得环 \(\mathbf{B}_{\lambda}\) 是 Noether 者组成）等于 \(\Lambda\) 。

因此我们有 \(\omega\in\Lambda'\) ，由此得到断言 \(a\) 和 \(b\) )。当 A 是 Noether 环时，我们有 \(\omega\in\Lambda''\) ，因而 B = B \(_{\omega}\) 是 Noether 的。

断言 c) 由 a)，b) 以及 III，§ 5，第 2 目的定理 1 得出。设 A（从而 B）是 Noether 的；由于我们有

\[
\left[ \mathfrak {m} _ {\mathrm{B}} ^ {n} / \mathfrak {m} _ {\mathrm{B}} ^ {n + 1}: \kappa_ {\mathrm{B}} \right] = \left[ \mathfrak {m} _ {\mathrm{A}} ^ {n} / \mathfrak {m} _ {\mathrm{A}} ^ {n + 1}: \kappa_ {\mathrm{A}} \right]
\]

对每个 \(n \in \mathbf{N}\) 成立，故 A 与 B 的 Hilbert-Samuel 级数相等。

推论. --- 假设 A 是 Noether 的。

\begin{enumerate}
\def\labelenumi{\alph{enumi})}
\item
  我们有 \(\dim (\mathbf{A}) = \dim (\mathbf{B})\)
\item
  假设 A 是正则的，并设 \((x_{1},\ldots ,x_{n})\) 是 A 的一个坐标系。则 B 是正则的，且序列 \((x_{1}1_{\mathrm{B}},\dots,x_{n}1_{\mathrm{B}})\) 是 B 的一个坐标系。
\end{enumerate}

这由 VIII，§ 5，第 1 目的命题 1 得出。

命题 3. --- 设 A，B，C 是三个局部环，\(u:A\to B\) ，\(v:B\to C\) 是两个膨胀。则 \(v\circ u\) 是一个膨胀。

设 \((\mathrm{B}_{\lambda})_{\lambda\in\Lambda}\) 与 \((\mathrm{C}_{\mu})_{\mu\in\mathbf{M}}\) 分别是 B 的子 A-代数族和 C 的子 B-代数族，具有定义 2 的性质 a)，b)，c)。在 \(\Lambda\) 与 M 的和 N 上，考虑这样的序关系：它在 \(\Lambda\) 和 M 上诱导给定的序，并且关系 \(\lambda<\mu\) 对 \(\lambda\in\Lambda\) ，\(\mu\in M\) 成立。这是一个良序关系。对 \(\lambda\in\Lambda\subset N\) ，令 \(C_{\lambda}=v(\mathrm{B}_{\lambda})\) 。则族 \((\mathrm{C}_{\nu})_{\nu\in\mathbf{N}}\) 满足关于 A-代数 C 的定义 1 的条件 a)，b)，c)。

定理 1. --- 设 \(f:A\to A'\) 是局部环之间的一个满射局部同态，B' 是 A' 的一个膨胀。则存在 A 的一个膨胀 B 以及从 B \(\otimes_{A}A'\) 到 B' 上的一个 A-代数同构。

\begin{enumerate}
\def\labelenumi{\Alph{enumi})}
\tightlist
\item
  假设 B' 是 A' 的一个初等膨胀。我们区分两种情形：1) 若 B' 在 A' 上有限，选取一个 A'-代数同构 \(\varphi: A'[X]/(P') \to B'\) ，其中 P' ∈ A'{[}X{]} 是一个首一多项式，其像在 \(\kappa_{A'}[X]\) 中不可约。选取一个首一多项式 P ∈ A{[}X{]} ，它在 A'{[}X{]} 中的像为 P'。它对 A 的极大理想取模必然不可约。令 B = A{[}X{]}/(P)。A-代数 B 是 A 的一个初等膨胀，且 \(\varphi\) 诱导从 \(B \otimes_{A} A'\) 到 B' 上的一个 A-代数同构。
\end{enumerate}

\begin{enumerate}
\def\labelenumi{\arabic{enumi})}
\setcounter{enumi}{1}
\tightlist
\item
  若 B' 在 A' 上不有限，选取一个 A'-代数同构 \(\psi : A']X[\to B'\) 。令 \(B = A]X[\) 。A-代数 B 是 A 的一个初等膨胀，且 B \(\otimes_{A} A'\) 典范地同构于 A'{]}X{[}。于是 \(\psi\) 诱导从 B \(\otimes_{A} A'\) 到 B' 上的一个 A-代数同构。
\end{enumerate}

\begin{enumerate}
\def\labelenumi{\Alph{enumi})}
\setcounter{enumi}{1}
\tightlist
\item
  现在过渡到一般情形。设 \((\mathbf{B}_{\lambda}^{\prime})_{\lambda\in\Lambda}\) 是由 \(A^{\prime}\) -代数组成的一个族，这些代数是 \(B^{\prime}\) 的子代数，它相对于 \(A^{\prime}\) 和 \(B^{\prime}\) 具有定义 2 的性质 a)，b)，c)。我们将用超限归纳定义一个归纳系统 \((\widetilde{\mathbf{B}}_{\lambda}, i_{\mu\lambda})\) ，它关于 \(\Lambda\) ，由局部环和单射的局部同态组成，以及同构 \(u_{\lambda}: \widetilde{B}_{\lambda} \otimes_{A} A^{\prime} \to B_{\lambda}^{\prime}\) ，使得对 \(\lambda \leqslant \mu\) ，\(u_{\mu} \circ (i_{\mu\lambda} \otimes \operatorname{Id}_{A^{\prime}}) \circ u_{\lambda}^{-1}\) 是 \(B_{\lambda}^{\prime}\) 到 \(B_{\mu}^{\prime}\) 中的典范单射。
\end{enumerate}

若 \(\alpha\) 是 \(\Lambda\) 的最小元，我们令 \(\tilde{\mathbf{B}}_{\alpha} = \mathbf{A}\) ，\(i_{\alpha \alpha} = \mathrm{Id}_{\mathbf{A}}\) ，并取 \(u_{\alpha}\) 为典范同构 \(\mathbf{A} \otimes_{\mathbf{A}} \mathbf{A}' \to \mathbf{A}'\) 。

设 \(v \in \Lambda\) ，并假设 \(\tilde{\mathbf{B}}_{\lambda}\) ，\(u_{\lambda}\) 和 \(i_{\mu \lambda}\) 当 \(\lambda \leqslant \mu < v\) 时均已构造。设 \(S_v\) 是由元素 \(\varepsilon\) 组成的集合，这些元素属于 \(\Lambda\) 且满足 \(\varepsilon < v\) 。若 \(S_v\) 没有最大元，则取 \(\tilde{\mathbf{B}}_v\) 为诸 \(\tilde{\mathbf{B}}_{\lambda}\) （\(\lambda \in S_v\) ）的归纳极限，取 \(u_v\) 为合成同构 \(\tilde{\mathbf{B}}_v \otimes_A A' \to \varinjlim (\tilde{\mathbf{B}}_{\lambda} \otimes_A A') \to \varinjlim B'_\lambda \to B'_v\) ，而取 \(i_{v\lambda}\) （当 \(\lambda \in S_v\) 时）为从 \(\tilde{\mathbf{B}}_{\lambda}\) 到 \(\tilde{\mathbf{B}}_v\) 的典范映射。若 \(S_v\) 有最大元 \(\mu\) ，则 \(B'_v\) 是 \(B'_\mu\) 的一个初等膨胀。由 A)，存在一个初等膨胀 \(i_{\nu \mu}:\tilde{\mathbf{B}}_{\mu}\to \tilde{\mathbf{B}}_{\nu}\) 以及一个 \(\tilde{\mathbf{B}}_{\mu}\) -代数同构，它把 \(\tilde{\mathbf{B}}_{\nu}\otimes_{\tilde{\mathbf{B}}_{\mu}}\mathbf{B}_{\mu}^{\prime}\) 映到 \(\mathbf{B}_{\nu}^{\prime}\) 上。取 \(u_{\nu}\) 为合成而得的 A-代数同构

\[
\widetilde {\mathbf {B}} _ {\nu} \otimes_ {\mathrm{A}} \mathrm{A} ^ {\prime} \rightarrow \widetilde {\mathbf {B}} _ {\nu} \otimes_ {\widetilde {\mathbf {B}} _ {\mu}} (\widetilde {\mathbf {B}} _ {\mu} \otimes_ {\mathrm{A}} \mathrm{A} ^ {\prime}) \rightarrow \widetilde {\mathbf {B}} _ {\nu} \otimes_ {\widetilde {\mathbf {B}} _ {\mu}} \mathrm{B} _ {\mu} ^ {\prime} \rightarrow \mathrm{B} _ {\nu} ^ {\prime}
\]

而取 \(i_{\nu \lambda}\) （当 \(\lambda \in S_{v}\) 时）为同态 \(i_{\nu \mu} \circ i_{\mu \lambda}\) 。

然后令 \(\mathbf{B} = \tilde{\mathbf{B}}_{\omega}\) ，且对所有 \(\lambda \in \Lambda\) ，用 \(\mathbf{B}_{\lambda}\) 表示 \(\tilde{\mathbf{B}}_{\lambda}\) 在典范单射 \(\tilde{\mathbf{B}}_{\lambda} \to \mathbf{B}\) 下的像。族 \((\mathbf{B}_{\lambda})_{\lambda \in \Lambda}\) 满足定义 2 的条件 \(a)\) ，\(b)\) ，\(c)\) ，且 \(\mathbf{B}\) 是 A 的一个膨胀。另一方面，同态 \(u_{\omega}\) 是从 \(\mathbf{B} \otimes_{A} \mathbf{A}'\) 到 \(\mathbf{B}'\) 上的一个 A'-同构。

推论. --- 设 A 是一个局部环，K 是其剩余域 \(\kappa_{\mathrm{A}}\) 的一个扩张。则存在一个局部环 B 以及一个膨胀 \(\mathbf{A} \to \mathbf{B}\) ，使得 \(\kappa_{\mathrm{A}}\) -代数 \(\kappa_{\mathrm{B}}\) 同构于 K。

事实上，同态 \(\kappa_{\mathbf{A}} \to \mathbf{K}\) 是一个膨胀（例 1）。取 \(\mathbf{A}' = \kappa_{\mathbf{A}}\) 和 \(\mathbf{B}' = \mathbf{K}\) 应用定理 1，我们得到 A 的一个膨胀 B 的存在性，以及从 \(\mathbf{B} / \mathfrak{m}_{\mathbf{A}}\mathbf{B}\) 到 K 上的一个 A-同构，由此得到本推论。

\refstepcounter{appendixitem}
